% Validation Framework
\section{Validation Framework}
\label{sec:validation}

SPOQ enforces quality through two \textit{validation gates}, structured
checkpoints where work is scored against explicit criteria before proceeding.
The first gate, \textit{epic validation}, occurs before execution and assesses
plan quality. The second, \textit{agent validation}, occurs after execution
and assesses code quality. Each gate applies 10 metrics, scored 0--100, with
a 95\% aggregate threshold for passing.

\subsection{Planning Validation: 10 Metrics}
\label{subsec:planning-validation}

Epic validation assesses whether the plan is sufficiently clear, complete,
and well-structured to enable successful execution. Table~\ref{tab:planning-metrics}
summarizes the metrics.

\begin{table}[htbp]
\centering
\caption{Epic Validation Metrics}
\label{tab:planning-metrics}
\small
\begin{tabular}{lp{6cm}l}
\toprule
\textbf{Metric} & \textbf{Question} & \textbf{Threshold} \\
\midrule
Vision Clarity (VC) & Is the epic's goal clearly scoped? & $\geq$90 \\
Architecture Quality (AQ) & Is the architecture diagram complete? & $\geq$90 \\
Task Decomposition (TD) & Are tasks atomic and independent? & $\geq$90 \\
Dependency Graph (DG) & Are dependencies explicit and acyclic? & $\geq$90 \\
Coverage Completeness (CC) & Do tasks fully cover all success criteria? & $\geq$90 \\
Phase Ordering (PO) & Do waves maximize parallelism? & $\geq$90 \\
Scope Coherence (SC) & Do all tasks contribute to the epic goal? & $\geq$90 \\
Success Criteria Quality (SQ) & Are criteria SMART? & $\geq$90 \\
Risk Identification (RI) & Are blockers and risks mitigated? & $\geq$90 \\
Integration Strategy (IS) & Is it clear how tasks merge and verify? & $\geq$90 \\
\bottomrule
\end{tabular}
\end{table}

\paragraph{Pass Criteria.}
An epic passes validation if:
\begin{equation}
\frac{1}{10}\sum_{i=1}^{10} M_i \geq 95 \quad \land \quad \min_i M_i \geq 90
\end{equation}

This dual requirement ensures both high average quality and no critically
weak dimensions. A perfect score on 9 metrics cannot compensate for a
failing score on the 10th.

\paragraph{Threshold Rationale.}
The 95/90 planning threshold reflects the cascading cost of planning errors:
a poorly-structured task creates execution problems that propagate to dependent
tasks. We observed that plans scoring below 90 on any metric required rework
in over 50\% of cases, while those scoring above 95 overall experienced less
than 10\% rework. The minimum threshold of 90 prevents ``averaging out'' a
critically weak dimension. Appendix~\ref{subsec:threshold-rationale}
provides additional design rationale.

\paragraph{Rationale: Validate Early.}
Planning mistakes are cheap to fix but expensive to execute. By enforcing
rigorous validation before spawning agents, SPOQ prevents wasted computation
on fundamentally flawed plans.

\newpage
\subsection{Code Validation: 10 Metrics}
\label{subsec:code-validation}

Agent validation assesses whether completed work meets quality standards.
Table~\ref{tab:code-metrics} summarizes the metrics.

\begin{table}[htbp]
\centering
\caption{Agent Validation Metrics}
\label{tab:code-metrics}
\small
\begin{tabular}{lp{6cm}l}
\toprule
\textbf{Metric} & \textbf{Question} & \textbf{Threshold} \\
\midrule
Syntactic Correctness (SC) & Does the code compile without errors? & $\geq$80 \\
Test Existence (TE) & Does new code have corresponding unit tests? & $\geq$80 \\
Test Pass Rate (TP) & Do all tests pass? & $\geq$80 \\
Requirements Fidelity (RF) & Does implementation match task specification? & $\geq$80 \\
SOLID Adherence (SA) & Does code follow SOLID design principles (Single responsibility, Open/closed, etc.)? & $\geq$80 \\
Security (SE) & Free from OWASP Top 10 vulnerabilities (injection, auth flaws, etc.)? & $\geq$80 \\
Error Handling (EH) & Does code handle failures gracefully? & $\geq$80 \\
Scalability (SL) & Will this code scale appropriately? & $\geq$80 \\
Code Clarity (CC) & Is code readable and self-documenting? & $\geq$80 \\
Completeness (CO) & Is work fully finished (no TODOs/stubs)? & $\geq$80 \\
\bottomrule
\end{tabular}
\end{table}

\paragraph{Pass Criteria.}
A task passes validation if:
\begin{equation}
\frac{1}{10}\sum_{i=1}^{10} M_i \geq 95 \quad \land \quad \min_i M_i \geq 80
\end{equation}

The per-metric floor is lower (80 vs 90) because code quality inherently
involves tradeoffs that planning does not.

\paragraph{Code Threshold Rationale.}
The 95/80 code threshold is more lenient than planning because code can be
iteratively improved post-delivery, and some metrics (e.g., security edge
cases) are inherently harder to achieve perfectly. Rework at the code level
is less expensive than re-planning an entire task decomposition; a failed
task can be re-executed with targeted feedback, whereas a flawed plan may
invalidate work across multiple dependent tasks.

\paragraph{Concise Feedback.}
On failure, the reviewer provides remediation guidance in $\leq$20 lines:
specific file:line references, concrete issues, and numbered action items.
This constraint respects the context budget of the orchestrator and forces
actionable specificity.

\newpage
\subsection{Validation Cascade}
\label{subsec:validation-cascade}

SPOQ applies a \textit{validation cascade}, a hierarchical check where the
overall epic validation incorporates individual task-level assessments.
Epic validation triggers task-level validation on each constituent task.
If any individual task scores $<$95, the epic's aggregate score is capped
at 85, forcing a FAIL verdict.

This prevents ``carrying'' weak tasks on the strength of strong planning.
Every task must be individually well-specified to pass epic validation.

\begin{figure}[htbp]
\centering
\begin{tikzpicture}[
    node distance=0.6cm,
    box/.style={rectangle, draw=spoqnavy, fill=spoqnavy!10, thick,
                minimum width=3cm, minimum height=0.6cm, align=center, font=\small},
    check/.style={circle, draw=spoqemerald, fill=spoqemerald!20, thick,
                  minimum size=0.4cm, inner sep=0pt},
    fail/.style={circle, draw=spoqburgundy, fill=spoqburgundy!20, thick,
                 minimum size=0.4cm, inner sep=0pt}
]

\node[box] (epic) {Epic Validation};
\node[box, below=of epic] (t1) {Task 01 Validation};
\node[box, below=of t1] (t2) {Task 02 Validation};
\node[box, below=of t2] (t3) {Task 03 Validation};
\node[below=0.3cm of t3, font=\small] (dots) {$\vdots$};

% Scores
\node[right=0.8cm of t1] (s1) {97 \checkmark};
\node[right=0.8cm of t2] (s2) {92 \checkmark};
\node[right=0.8cm of t3, text=spoqburgundy] (s3) {84 $\times$};

% Cascade effect
\draw[-Stealth, dashed, spoqburgundy] (s3.west) -- ++(-0.3,0) -- ++(0,2.5) -- (epic.east);
\node[right=1.5cm of epic, text=spoqburgundy, font=\small] {Capped at 85};

\end{tikzpicture}
\caption{Validation cascade: a single task scoring $<$95 caps the epic's
overall score, triggering failure.}
\label{fig:cascade}
\end{figure}
